\documentclass[11pt]{article}
\usepackage[margin=2.5cm]{geometry}
\usepackage{amsmath}
\usepackage{parskip}

\title{Mixed models and the random intercept, simply}
\date{}

\begin{document}
\maketitle

\section*{The problem: rows that belong together}

In this thesis every person contributes nine rows, one per segment. Two
segments of the same person are not independent. They share the same blood,
blood pressure, genes and lifestyle, much of which the model does not
measure. If one segment of a person developed plaque, the other eight are
more likely to have done so too. A model that treats all rows as independent
ignores this and believes it has more information than it has: its standard
errors come out too small and its tests too optimistic.

\section*{The idea: every person gets their own baseline}

A random intercept gives each person $i$ their own shift $b_i$ of the risk
scale, added to the usual linear predictor of segment $s$,
\[
  \eta_{is} = \beta_0 + \beta_1 x_{is,1} + \beta_2 x_{is,2} + \dots + b_i .
\]
A person with a positive $b_i$ has a higher risk in all nine segments than
the features alone predict; a person with a negative $b_i$ has a lower risk
in all nine. The shift is the same for every segment of that person, and this
is what ties the segments of one person together. The $b_i$ soak up the
unmeasured, person-wide part of the risk.

\section*{Why ``random''}

There are about a thousand people, and estimating a separate free intercept
for each would cost a thousand parameters, most of them from only nine
segments. Instead the model assumes that the shifts come from one common
distribution,
\[
  b_i \sim N(0, \sigma^2),
\]
and estimates only its spread $\sigma$. A large $\sigma$ means people differ a
lot in their overall plaque tendency beyond the measured features; $\sigma = 0$
means they do not, and the model reduces to an ordinary GLM. The individual
$b_i$ are not parameters of the model; they are averaged over.

\section*{Why ``mixed''}

The model mixes two kinds of effects. The \emph{fixed} effects, the $\beta$s,
are the same for everyone: the effect of curvature, taper or SCORE2. The
\emph{random} effects, the $b_i$, vary from person to person. A generalised
linear mixed model (GLMM) is a GLM with such random effects added.

\section*{What it means for the coefficients}

Because $b_i$ is held fixed when reading a coefficient, $e^{\beta}$ compares
two segments with the same person-level shift: it is a \emph{conditional}
effect. Averaged over all people, the effect on the population risk is
somewhat smaller. When a new person's risk is predicted, their own $b_i$ is
unknown, so the model averages the risk over all possible values of $b$;
the thesis model does this when it pools the nine segment risks into one
person risk.

\end{document}
